Throughout, a nondegenerate continuous symmetric form on a bundle with a bundle metric has a continuous self-adjoint invertible map; its spectral projections onto the positive and negative parts are continuous and of locally constant rank, no eigenvalue crossing zero, so the positive and negative eigenspaces form vector subbundles. In particular the signature is locally constant, hence constant on a connected base: the first step of the proof of Theorem \claimref{sheaf.conservation-degeneracy-absolute-form}, which uses no tameness.

\emph{(1).} Let $h$ be a strict completion standard outside $U$, nondegenerate at every point of $M$. At $x, y \in M \setminus U$ its signature is $\operatorname{sig} G(x) + \operatorname{sig} K$ and $\operatorname{sig} G(y) + \operatorname{sig} K$ with the same $\operatorname{sig} K$; these differ if $\operatorname{sig} G(x) \neq \operatorname{sig} G(y)$, against constancy. So the set $D$ where $h$ degenerates is nonempty; it is closed (the zero set of a local Gram determinant) and lies in $U$, since on $M \setminus U$ both summands of $G \perp K$ are nondegenerate ($\Sig \subset U$). Pulling the bundle and the form back to a path from $x$ to $y$ and applying the constancy on the interval shows that the path meets $D$. The stored theorem is the case of a collar and the values $(p,q)$, $(p-1,q+1)$. For the place of the degeneration: $G \perp K$ with no deformation is a strict completion and degenerates exactly on $\Sig$; and $\det\begin{psmallmatrix} x & \beta \\ \beta & \pm 1\end{psmallmatrix} = \pm x - \beta^{2}$ vanishes exactly where $x = \pm\beta^{2}$, which for $\beta$ positive at $0$ and zero outside $U$ is a point with $x > 0$, respectively $x < 0$.

\emph{(2).} A continuous $\beta$ exists by Urysohn's lemma, $\Sig$ and $M \setminus U$ being disjoint closed sets of a manifold ($\beta = 0$ if $\Sig$ is empty). On $V \oplus V$ with the metric $A \oplus A$, $\varepsilon_{1}(a \oplus a') = a \oplus (-a')$ and $\varepsilon_{2}(a \oplus a') = a' \oplus a$ satisfy $\varepsilon_{1}^{2} = \varepsilon_{2}^{2} = 1$, $\varepsilon_{1}\varepsilon_{2} = -\varepsilon_{2}\varepsilon_{1}$, both are self-adjoint and both commute with $g$ acting on each copy. The map of $h_{D}$ is $\hat h = g\varepsilon_{1} + \beta\varepsilon_{2}$, and $\hat h^{2} = g^{2} + \beta^{2} + g\beta(\varepsilon_{1}\varepsilon_{2} + \varepsilon_{2}\varepsilon_{1}) = g^{2} + \beta^{2}$, positive definite wherever $g$ is invertible or $\beta \neq 0$, that is everywhere. So $\hat h$ is invertible; and $\varepsilon_{1}\varepsilon_{2}$ anticommutes with $\hat h$ and is invertible, so it exchanges the eigenspaces of opposite eigenvalues and the signature is $(n,n)$. Outside $U$, $\beta = 0$ and $h_{D} = G \perp (-G)$. This is the form of clause (2) of Theorem \claimref{synthesis.wall-class-one-rung-obstruction}; tameness is not used.

\emph{(3).} If $\operatorname{sig} G$ is constant on $M \setminus U$, the double's $K = -G$ on the second copy is nondegenerate there with one signature: the double is strict. The converse is (1).

\emph{(4).} By the opening paragraph $h_{\pm}$ are bundles over $M$. On $M \setminus U$ take the auxiliary metric to be a sum of metrics on $V$ and on the complement; then the map of $h = G \perp K$ is the sum of the maps of $G$ and $K$, and its positive and negative eigenbundles are $G_{\pm} \oplus K_{\pm}$. The eigenbundles for two metrics are isomorphic: each is a maximal positive (or negative) subbundle, and the projection of one maximal positive subbundle to another along a maximal negative one is injective, a positive vector in a negative subbundle being zero. So these extend over $M$. Over a contractible base every vector bundle is trivial; then the ranks are constant, which is the constancy of $\operatorname{sig} G + \operatorname{sig} K$, and $0 = w_{1}(G_{\pm} \oplus K_{\pm}) = w_{1}(G_{\pm}) + w_{1}(K_{\pm})$ in the cohomology of $M \setminus U$ modulo two.

\emph{(5).} The form has trace $0$ and determinant $-(x^{2} + y^{2})$: signature $(1,1)$ off the origin and zero at it. At the point $r(\cos\theta, \sin\theta)$, with $c = \cos\tfrac\theta2$, $s = \sin\tfrac\theta2$, the matrix is $r\begin{psmallmatrix} c^{2} - s^{2} & 2cs \\ 2cs & s^{2} - c^{2}\end{psmallmatrix}$, and it sends $(c,s)$ to $r(c,s)$ and $(-s,c)$ to $-r(-s,c)$. So $G_{+}$ is spanned by $(c,s)$, a vector that depends continuously on $\theta$ and at $\theta = 2\pi$ is the negative of its value at $\theta = 0$: around a circle centred at the origin the line $G_{+}$ has no continuous nonvanishing section, and neither has $G_{-}$. Let $h$ be a model completion standard outside a bounded $U$, nondegenerate everywhere, and let $C$ be a circle centred at the origin outside $U$. By (4), $h_{+}$ is a bundle over the plane, hence trivial, hence orientable, and so is its restriction to $C$. But that restriction is $G_{+} \oplus (K_{0})_{+}$ with $(K_{0})_{+}$ a fixed subspace: carrying the frame $\bigl((c,s), k_{1}, \dots, k_{a}\bigr)$ once around $C$ returns it with its first vector reversed, so the bundle is not orientable. Contradiction. For the double, $K = -G$ and $K_{+} = G_{-}$, spanned by $(-s,c)$, which is reversed likewise: $G_{+} \oplus K_{+}$ is the whole of $V$ on $C$, trivial.

\emph{(6).} (a) $x^{2} > 0$ for $x \neq 0$, corank one at $0$, vanishing to second order, so the section of (T3) vanishes; the determinant $-x^{2} - \beta^{2}$ is negative where $x \neq 0$ or $\beta \neq 0$, and the form is $x^{2} \perp (-1)$ where $\beta = 0$. (b) The eigenvalues are $x$ and $-x$: signature $(1,1)$ off $0$ and corank two at $0$. For $h = \begin{psmallmatrix} G & \beta C \\ \beta C^{\mathsf T} & K_{0}\end{psmallmatrix}$ one has $\det h = \det K_{0}\cdot\det(G - \beta^{2}CK_{0}^{-1}C^{\mathsf T})$; with $C = \begin{psmallmatrix} 1 & 1 \\ 1 & -1\end{psmallmatrix}$ and $K_{0} = \operatorname{diag}(1,-1)$, $CK_{0}^{-1}C^{\mathsf T} = \begin{psmallmatrix} 0 & 2 \\ 2 & 0\end{psmallmatrix}$ and $\det h = -\bigl(-x^{2} - 4\beta^{4}\bigr) = x^{2} + 4\beta^{4} > 0$ where $x \neq 0$ or $\beta \neq 0$; the signature is $(2,2)$, its value where $\beta = 0$. (c) $x^{2} + y^{2} > 0$ off the origin; as in (a). In each case $\beta$ is positive on $\Sig$ and zero outside $U$, and the complement and its form outside $U$ are fixed.

\emph{(7).} For $x \neq 0$ the form is definite, of the sign of $x$; (1) applies with any $U \ni 0$ whose complement meets both half-lines. The residues are $4$ and $-4 \equiv 4$, and the two algebras are isomorphic by the classification (\claimref{std-algebra.periodicity}). In the tame case $(p,q) \to (p-1,q+1)$ changes $p - q$ by $2$.

\emph{(8).} $g(x)^{2} = \sum_{i}x_{i}^{2}E_{i}^{2} + \sum_{i<j}x_{i}x_{j}(E_{i}E_{j} + E_{j}E_{i}) = Q(x)$, so $g(x)$ is invertible exactly for $x \neq 0$, and $g(0) = 0$ has corank $d$. If $E_{m+1}$ is symmetric, of square one and anticommutes with each $E_{i}$, then $(g + \beta E_{m+1})^{2} = Q(x) + \beta^{2}$, positive where $x \neq 0$ or $\beta \neq 0$; the form is $G$ where $\beta = 0$, so it is a model completion with $W_{0} = 0$. Symmetric generators exist on every module for a suitable inner product (average one over the finite group the generators generate; an orthogonal map of square one is symmetric). The least module of $\Cl(m,0)$ has dimension $1, 2, 4, 8, 8, 16, 16, 16, 16$ for $m = 1, \dots, 9$ and sixteen times that eight steps later, from the table of \claimref{std-algebra.periodicity} ($\R \oplus \R$, $\Mat_{2}(\R)$, $\Mat_{2}(\C)$, $\Mat_{2}(\HH)$, $\Mat_{2}(\HH) \oplus \Mat_{2}(\HH)$, $\Mat_{4}(\HH)$, $\Mat_{8}(\C)$, $\Mat_{16}(\R)$, $\Mat_{16}(\R) \oplus \Mat_{16}(\R)$). Where the dimensions for $m$ and $m + 1$ agree, $m \equiv 0, 4, 6, 7$, the algebra $\Cl(m,0)$ is simple, so its least module is unique and is the restriction of a least module of $\Cl(m+1,0)$, which supplies a map $A$ of square one anticommuting with $E_{1}, \dots, E_{m}$. To have it symmetric for the inner product in which the $E_{i}$ are symmetric, put $S = A + A^{\mathsf T}$: it is symmetric and anticommutes with each $E_{i}$, its kernel is invariant under the $E_{i}$ and so is zero or everything on the irreducible module, and it is not zero, since $A^{\mathsf T} = -A$ with $A^{2} = 1$ would give $A^{\mathsf T}A = -1$; $S^{2}$ is symmetric, positive and commutes with the $E_{i}$, and $E_{m+1} = S\,(S^{2})^{-1/2}$ is symmetric, of square one and anticommutes with each $E_{i}$. Conversely, where the dimension for $m + 1$ is larger no such map exists on the least module, which would otherwise be a module of $\Cl(m+1,0)$ of too small a dimension. For $m = 1, 2, 3$ the space of symmetric maps anticommuting with the generators on the least module is zero (computed), and for $m = 4$ it is a line spanned by a map of square one (computed). For $m = 1$ the family is $\pm x$, of signatures $(1,0)$ and $(0,1)$ on the two sides: (1), for $U$ bounded, so that its complement meets both sides. For $m = 2$ on the least module the generators can be taken as $\operatorname{diag}(1,-1)$ and the exchange, and $g$ is the family of (5).

\emph{(9).} $V_{x}$ with $G_{x}$ is a quadratic space with radical of dimension $r$ inside the nondegenerate $(E_{x}, h_{x})$; by Fact \claimref{odd-extensions.imported-minimal-completion-degenerate-algebra-rule} the dimension is at least $(n - r) + 2r$, and, with $N$ a complement of the radical in $V_{x}$, of signature $(p,q)$, the orthogonal of $N$ in $E_{x}$ is nondegenerate and contains the radical, a totally isotropic subspace of dimension $r$, so it has at least $r$ directions of each sign; adding those of $N$ gives $p + r$ and $q + r$.

\emph{Reading.} The conjecture's strict reading is (1) with ``one side'' and ``the other side'' for two exterior regions. Its failure as a general statement is (6); its true general form is (4); (7) and (5) are the two statements about the invariant. The pattern in $m$ is (8); the unproved half is marked as cited.
